\section{Traffic Analysis Results}

This section presents the five required Wireshark captures. Each subsection
below has the exact setup already filled in, since that part does not depend
on the actual capture, only on what to insert once the capture exists.

\subsection{Protocol Distribution}

Capture at least 60 seconds with the \texttt{balanced} profile running all
four generators. Export Statistics $\rightarrow$ Protocol Hierarchy as a
screenshot.

\todoblock{INSERT}{Screenshot of the Protocol Hierarchy and the percentage
breakdown of TCP vs.\ UDP, and within TCP the HTTP/2 (port 8080), MQTT (port
1883) and raw TCP shares, plus the UDP-only QUIC share (port 4433).}

\subsection{Temporal Analysis}

Capture 120 seconds with the \texttt{http2\_heavy} profile, which includes a
periodic-burst phase (burst rate 400 for 5 seconds every 30 seconds). Export
Statistics $\rightarrow$ I/O Graph with one line per protocol filter.

\todoblock{INSERT}{I/O graph screenshot with the burst peaks marked, plus a
second capture of a warmup period and a Random/Poisson phase for comparison
against the constant and burst shapes.}

\subsection{Behavioral Fingerprinting}

Capture two 30-second windows of the TCP/UDP generator alone: one in
\texttt{mode=normal}, one in \texttt{mode=stealth}. Optionally extend this with
the \texttt{balanced\_with\_stealth} phase, which runs stealth mode together
with the random pattern on HTTP/2, QUIC and MQTT at the same time.

\todoblock{INSERT}{Packet-length histograms and inter-arrival-time
distributions for both modes, with an explicit statement of which
characteristics are distinguishable (e.g.\ the fixed 100\,ms spacing in normal
mode) and which overlap with real client traffic.}

\subsection{Failure Visibility}

Capture across a \texttt{docker stop target-http2} event during normal
operation. Filter on \texttt{tcp.flags.reset==1} for the TCP RST, or look for
repeated SYNs to port 1883 if stopping \texttt{mosquitto} instead.

\todoblock{INSERT}{Measured time from the stop command to the first visible
failure indicator, and which protocol-specific signal appeared. See Section VI
for the deeper discussion of this capture together with the Adaptive Control
reaction.}

\subsection{Multi-System Deployment}

Capture on the physical network interface of either lab machine (not
\texttt{docker0}/\texttt{br-*}) while \texttt{docker-compose.targets.yml} runs
on Machine B and \texttt{docker-compose.generators.yml} runs on Machine A with
\texttt{TARGET\_B\_IP} pointing at Machine B.

\todoblock{INSERT}{Capture screenshot plus measured latency per protocol,
compared against an equivalent single-machine run. See Section VII for the
full comparison.}
